\setcounter{chapter}{4}
\chapter{Lebesgue測度}

\paragraph{本章の目標}
有界な窓で切ったときにLebesgue外測度 $\lambda^*$ とLebesgue内測度 $\lambda_*$ が一致する集合を
Lebesgue可測集合と定める．
そのうえで，Lebesgue可測集合全体が $\sigma$-加法族をなし，
Lebesgue測度が可算加法的な測度になることを確認する．
最後に，内外正則性とCarathéodory条件による判定法を述べる．

\section{Lebesgue可測集合とLebesgue測度}

以下では，
\[
Q_R:=[-R,R]^d\qquad(R>0)
\]
と書く．

\begin{definition}[Lebesgue可測集合とLebesgue測度]\label{def:lebesgue-measurable-local}
集合 $A \subset \RR^d$ がLebesgue可測であるとは
\[
\lambda_*(A\cap Q_R)=\lambda^*(A\cap Q_R)
\qquad (R>0)
\]
が成り立つことをいう．
Lebesgue可測集合全体を
\[
\calL_d:=\{A\subset\RR^d\mid A \text{ はLebesgue可測}\}
\]
と書く．
$A\in\calL_d$ のとき，
\[
\lambda(A):=\lambda^*(A)
\]
を $A$ のLebesgue測度という．
\end{definition}

% \begin{lemma}[有界集合の可測性]\label{lem:bounded-measurable-iff-inner-outer}
% 有界集合 $A\subset\RR^d$ について，次は同値である．
% \begin{enumerate}
%     \item $A$ はLebesgue可測である
%     \item $\lambda_*(A)=\lambda^*(A)$
% \end{enumerate}
% また，$A$ が有界ならば 
% %\[
% $\lambda_*(A)<\infty$
% %\]
% である．
% \end{lemma}
% \begin{proof}
% $A$ は有界なので，ある $R_0>0$ が存在して
% \[
% A\subset Q_{R_0}=[-R_0,R_0]^d
% \]
% となる．

% まず $A$ がLebesgue可測であるとする．
% 定義より
% \[
% \lambda_*(A\cap Q_{R_0})=\lambda^*(A\cap Q_{R_0})
% \]
% である．
% $A\cap Q_{R_0}=A$ だから
% \[
% \lambda_*(A)=\lambda^*(A)
% \]
% を得る．

% 逆に $\lambda_*(A)=\lambda^*(A)$ とする．
% 区間による外測度の分解（\cref{cor:interval-outer-measure-split}）を
% $A\subset Q_{R_0}$ に適用すると
% \[
% \lambda^*(Q_{R_0})
% =
% \lambda^*(A)+\lambda^*(Q_{R_0}\setminus A)
% \]
% である．

% 任意の $R>0$ に対し，
% \[
% A_R:=A\cap Q_R
% \]
% とおく．
% $A_R$ が内外一致することを示せばよい．
% $R>R_0$ ならば $A_R=A$ だから，仮定より内外一致する．
% 以下では $R\le R_0$ とする．

% コンパクト集合による外測度の分解
% （\cref{lem:compact-splits-outer-measure}）を
% コンパクト集合 $Q_R$ と集合 $A,\,Q_{R_0}\setminus A,\,Q_{R_0}$ にそれぞれ適用すると，
% \[
% \lambda^*(A)=\lambda^*(A_R)+\lambda^*(A\setminus Q_R),
% \]
% \[
% \lambda^*(Q_{R_0}\setminus A)
% =
% \lambda^*(Q_R\setminus A_R)
% +
% \lambda^*((Q_{R_0}\setminus A)\setminus Q_R),
% \]
% \[
% \lambda^*(Q_{R_0})
% =
% \lambda^*(Q_R)+\lambda^*(Q_{R_0}\setminus Q_R)
% \]
% である．ここで
% \[
% Q_{R_0}\setminus Q_R
% =
% (A\setminus Q_R)\cup((Q_{R_0}\setminus A)\setminus Q_R)
% \]
% だから，外測度の劣加法性より
% \[
% \lambda^*(Q_{R_0}\setminus Q_R)
% \le
% \lambda^*(A\setminus Q_R)
% +
% \lambda^*((Q_{R_0}\setminus A)\setminus Q_R)
% \]
% である．
% すべて $Q_{R_0}$ の部分集合なので外測度は有限であり，上の等式を比較すると
% \[
% \lambda^*(A_R)+\lambda^*(Q_R\setminus A_R)
% \le
% \lambda^*(Q_R)
% \]
% を得る．逆向きの不等式は外測度の劣加法性から従うので
% \[
% \lambda^*(Q_R)
% =
% \lambda^*(A_R)+\lambda^*(Q_R\setminus A_R)
% \]
% である．
% 最後に，第3章で示した区間の外測度の計算より $\lambda^*(Q_R)=|Q_R|$ であるから，
% 区間を用いた内測度表示を $A_R\subset Q_R$ に適用すれば
% \[
% \lambda_*(A_R)
% =
% \lambda^*(Q_R)-\lambda^*(Q_R\setminus A_R)
% =
% \lambda^*(A_R)
% \]
% を得る．
% よって $A_R=A\cap Q_R$ は内外一致する．
% $R>0$ は任意だから，$A$ はLebesgue可測である．

% 最後に，$A\subset Q_{R_0}$ と
% Lebesgue外測度の単調性（\cref{thm:lebesgue-outer-measure}）より
% \[
% \lambda_*(A)\le \lambda^*(A)\le \lambda^*(Q_{R_0})<\infty
% \]
% である．
% \end{proof}

% \begin{corollary}[有界集合の区間判定]\label{cor:bounded-measurable-interval-criterion}
% 有界集合 $A\subset\RR^d$ と，$A\subset Q$ を満たす有界区間 $Q$ に対して，
% 次は同値である．
% \begin{enumerate}
%     \item $A$ はLebesgue可測である
%     \item $\lambda^*(A)+\lambda^*(Q\setminus A)=|Q|$
% \end{enumerate}
% \end{corollary}
% \begin{proof}
% 有界集合の可測性（\cref{lem:bounded-measurable-iff-inner-outer}）と
% 区間による外測度の分解（\cref{cor:interval-outer-measure-split}）より，
% $A$ がLebesgue可測であることは
% \[
% \lambda^*(A)+\lambda^*(Q\setminus A)=|Q|
% \]
% と同値である．
% \end{proof}

\begin{theorem}[有限外測度集合の可測性]\label{thm:finite-outer-measurable-iff-inner-outer}
$\lambda^*(A)<\infty$ とする．このとき次は同値である．
\begin{enumerate}
    \item $A$ はLebesgue可測である
    \item $\lambda_*(A)=\lambda^*(A)$
\end{enumerate}
\end{theorem}
\begin{proof}
まず $A$ がLebesgue可測であるとする．
$\eps>0$ を任意に取る．
Lebesgue外測度の定義より，$A$ を覆う区間列 $\{I_n\}_{n=1}^{\infty}$ で
\[
\sum_{n=1}^{\infty}|I_n|<\lambda^*(A)+\eps
\]
となるものが取れる．
ある $N$ について
\[
\sum_{n>N}|I_n|<\eps
\]
である．さらに $R>0$ を十分大きく取って，
空でない $I_1,\dots,I_N$ がすべて $Q_R$ に含まれるようにする．
すると
\[
A\setminus Q_R\subset\bigcup_{n>N} I_n
\]
だから
\[
\lambda^*(A\setminus Q_R)<\eps.
\]
$A$ はLebesgue可測なので，$A\cap Q_R$ は内外一致する．
したがって，コンパクト集合 $K\subset A\cap Q_R$ で
\[
\lambda^*(K)>\lambda^*(A\cap Q_R)-\eps
\]
となるものが取れる．外測度の劣加法性より
\[
\lambda^*(A)
\le
\lambda^*(A\cap Q_R)+\lambda^*(A\setminus Q_R)
<
\lambda^*(K)+2\eps
\le
\lambda_*(A)+2\eps.
\]
$\eps>0$ は任意なので，Lebesgue内測度の基本性質
（\cref{prop:lebesgue-inner-basic}）と合わせて
\[
\lambda_*(A)=\lambda^*(A)
\]
である．

逆に $\lambda_*(A)=\lambda^*(A)$ とする．
任意の $R>0$ に対し，
\[
A_R:=A\cap Q_R
\]
が内外一致することを示す．
$\eps>0$ を任意に取る．内測度の定義より，コンパクト集合 $K\subset A$ で
\[
\lambda^*(K)>\lambda^*(A)-\eps
\]
となるものが取れる．
コンパクト集合による外測度の分解（\cref{lem:compact-splits-outer-measure}）を
コンパクト集合 $Q_R$ と集合 $A,\,K$ にそれぞれ適用すると，
\[
\lambda^*(A)=\lambda^*(A_R)+\lambda^*(A\setminus Q_R),
\]
\[
\lambda^*(K)
=
\lambda^*(K\cap Q_R)+\lambda^*(K\setminus Q_R)
\]
である．また $K\setminus Q_R\subset A\setminus Q_R$ だから
\[
\lambda^*(K\cap Q_R)
>
\lambda^*(A)-\lambda^*(A\setminus Q_R)-\eps
=
\lambda^*(A_R)-\eps.
\]
$K\cap Q_R$ は $A_R$ に含まれるコンパクト集合なので，
$\lambda_*(A_R)>\lambda^*(A_R)-\eps$ である．
$\eps>0$ は任意だから
\[
\lambda_*(A_R)\ge \lambda^*(A_R)
\]
を得る．逆向きの不等式はLebesgue内測度の基本性質であるから，
$A_R$ は内外一致する．
$R>0$ は任意だから，$A$ はLebesgue可測である．
\end{proof}

\begin{remark}
一般に $A$ が有界ならば外測度有限 $\lambda^*(A)<\infty$ である．しかし，逆
\[
\lambda^*(A)<\infty \implies A \text{ は有界}
\]
は成り立たない．
例えば $\ZZ^d$ は非有界であるが，可算集合なので
\[
\lambda^*(\ZZ^d)=0
\]
である．

また，有限外測度集合の可測性
（\cref{thm:finite-outer-measurable-iff-inner-outer}）では
$\lambda^*(A)<\infty$ という仮定が本質的である．
後で示すLebesgue非可測集合 $V\subset[0,1]\subset\RR$ を用いて
\[
A:=V\cup[2,\infty)\subset\RR
\]
とおく．このとき任意の $N>2$ に対して $[2,N]\subset A$ だから
\[
\lambda_*(A)\ge \lambda^*([2,N])=N-2
\]
であり，$N\to\infty$ として $\lambda_*(A)=\infty$ である．
したがって $\lambda^*(A)=\infty$ でもあり，
\[
\lambda_*(A)=\lambda^*(A)=\infty
\]
が成り立つ．一方，
\[
A\cap[-1,1]=V
\]
であるから，$A$ はLebesgue可測ではない．
\end{remark}

% \begin{definition}[Lebesgue零集合]
% 集合 $N\subset\RR^d$ が
% \[
% \lambda^*(N)=0
% \]
% を満たすとき，$N$ をLebesgue零集合という．
% \end{definition}

\begin{theorem}[Lebesgue零集合]\label{thm:lebesgue-null-measurable}
Lebesgue零集合はLebesgue可測であり，Lebesgue測度は $0$ である．
\end{theorem}
\begin{proof}
任意の $R>0$ に対して，Lebesgue外測度の単調性
（\cref{thm:lebesgue-outer-measure}）より
\[
\lambda^*(N\cap[-R,R]^d)\le \lambda^*(N)=0
\]
である．
Lebesgue内測度の基本性質（\cref{prop:lebesgue-inner-basic}）より
\[
0\le \lambda_*(N\cap[-R,R]^d)
\le
\lambda^*(N\cap[-R,R]^d)
\]
である．
したがって
\[
\lambda_*(N\cap[-R,R]^d)
=
\lambda^*(N\cap[-R,R]^d)
=0
\]
だから，$N$ はLebesgue可測である．
また $\lambda(N)=\lambda^*(N)=0$ である．
\end{proof}

\begin{corollary}[可算集合]\label{cor:countable-lebesgue-null}
一点集合，可算集合，$\QQ^d$，$\ZZ^d$ はLebesgue可測であり，測度 $0$ をもつ．
%特に1次元では $\QQ$ と $\ZZ$ はLebesgue測度 $0$ である．
\end{corollary}
\begin{proof}
一点集合および可算集合のLebesgue外測度は $0$ である
（\cref{thm:countable-outer-measure-zero}）．
したがってLebesgue零集合の可測性（\cref{thm:lebesgue-null-measurable}）より従う．
\end{proof}

\begin{theorem}[Jordan可測集合]\label{thm:jordan-measurable-lebesgue}
Jordan可測集合 $A\subset\RR^d$ はLebesgue可測であり，
\[
\lambda(A)=m_J(A)
\]
が成り立つ．
\end{theorem}
\begin{proof}
    Jordan可測集合は有界なのでLebesgue外測度有限である（$\lambda^*(A) < \infty$）．よって$\lambda_*(A)=\lambda^*(A)$を示せばよい．
Jordan可測性より
\[
m_{J,*}(A)=m_J^*(A)=m_J(A)
\]
である．
第4章のJordan内測度との比較，Lebesgue内測度の基本性質，
第3章のJordan外測度との比較を順に用いると
\[
m_J(A)
=m_{J,*}(A)
\le \lambda_*(A)
\le \lambda^*(A)
\le m_J^*(A)
=m_J(A)
\]
となる．
したがってすべて等号であり，特に
\[
\lambda_*(A)=\lambda^*(A)=m_J(A)
\]
である．
よって $A$ はLebesgue可測であり，
\[
\lambda(A)=m_J(A)
\]
である．
\end{proof}
% \begin{proof}
% Jordan可測集合は有界なので，有界閉区間 $Q$ を
% \[
% A\subset Q
% \]
% となるように取る．
% $Q\setminus A$ もJordan可測であり，Jordan測度の有限加法性より
% \[
% m_J(A)+m_J(Q\setminus A)=m_J(Q)=|Q|
% \]
% である．
% 第3章で示したJordan測度とLebesgue外測度の一致より
% \[
% \lambda^*(A)+\lambda^*(Q\setminus A)=|Q|
% \]
% だから，有界集合の区間判定（\cref{cor:bounded-measurable-interval-criterion}）より
% $A$ はLebesgue可測である．
% さらに
% \[
% \lambda(A)=\lambda^*(A)=m_J(A)
% \]
% は，ふたたびJordan測度とLebesgue外測度の一致から従う．
% \end{proof}

\begin{example}
\[
\QQ\cap[0,1]
\]
は可算集合なのでLebesgue可測であり，
\[
\lambda_1(\QQ\cap[0,1])=0
\]
である．
一方，$[0,1]$ はJordan可測で $m_J([0,1])=1$ だから
\[
\lambda_1([0,1])=1
\]
である．
よって，後で示す完全加法性（\cref{thm:lebesgue-measure-countable-additivity}）から従う有限加法性より
\[
\lambda_1([0,1]\setminus\QQ)=1
\]
となる．
\end{example}

\section{Lebesgue可測集合の性質}

まず，有界な窓の中で使う有限外測度の判定法を用意する．

\begin{lemma}[有限外測度集合の近似判定法]\label{lem:finite-outer-approx-criterion}
$E\subset\RR^d$ が $\lambda^*(E)<\infty$ を満たすとする．
このとき次は同値である．
\begin{enumerate}
    \item $\lambda_*(E)=\lambda^*(E)$
    \item 任意の $\eps>0$ に対して，コンパクト集合 $K$ と開集合 $G$ が存在して
    \[
    K\subset E\subset G,
    \qquad
    \lambda^*(G\setminus K)<\eps
    \]
    を満たす
    \item 任意の $\eps>0$ に対して，開集合 $G\supset E$ が存在して
    \[
    \lambda^*(G\setminus E)<\eps
    \]
    を満たす
\end{enumerate}
\end{lemma}
\begin{proof}
(1) から (2) を示す．
$E$ が内外一致するとする．
内測度の定義より，コンパクト集合 $K\subset E$ で
\[
\lambda^*(K)>\lambda^*(E)-\frac{\eps}{2}
\]
を満たすものが取れる．
また開集合による外正則性（\cref{thm:lebesgue-outer-regularity-open}）より，開集合 $G\supset E$ で
\[
\lambda^*(G)<\lambda^*(E)+\frac{\eps}{2}
\]
を満たすものが取れる．
特に $\lambda^*(G)<\infty$ である．
開集合とコンパクト集合の分解公式（\cref{lem:finite-open-compact-split}）を $G$ と $K$ に適用すると
\[
\lambda^*(G)=\lambda^*(K)+\lambda^*(G\setminus K)
\]
である．
したがって
\[
\lambda^*(G\setminus K)
=\lambda^*(G)-\lambda^*(K)
<\eps
\]
を得る．

(2) から (3) は
\[
G\setminus E\subset G\setminus K
\]
より従う．

(3) から (1) を示す．
$\eps>0$ を任意に取る．
$G\supset E$ を
\[
\lambda^*(G\setminus E)<\frac{\eps}{3}
\]
となる開集合とする．
外測度の劣加法性より
\[
\lambda^*(G)\le \lambda^*(E)+\lambda^*(G\setminus E)<\infty
\]
である．
開集合による外正則性より，$G\setminus E$ を含む開集合 $U$ で
\[
\lambda^*(U)<\lambda^*(G\setminus E)+\frac{\eps}{3}
\]
となるものを取る．
また \cref{lem:finite-open-inner-compact} より，コンパクト集合 $F\subset G$ で
\[
\lambda^*(F)>\lambda^*(G)-\frac{\eps}{3}
\]
となるものを取る．

$K:=F\setminus U$ とおくと，$K$ は $E$ に含まれるコンパクト集合である．
さらに $F\subset K\cup U$ だから
\[
\lambda^*(K)
\ge
\lambda^*(F)-\lambda^*(U)
>
\lambda^*(G)-\lambda^*(G\setminus E)-\frac{2\eps}{3}
\ge
\lambda^*(E)-\eps.
\]
内測度の定義より $\lambda_*(E)\ge\lambda^*(K)$ だから，
$\lambda_*(E)>\lambda^*(E)-\eps$ である．
$\eps>0$ は任意だから
\[
\lambda_*(E)\ge \lambda^*(E)
\]
である．
逆向きの不等式はLebesgue内測度の基本性質（\cref{prop:lebesgue-inner-basic}）であるから，
\[
\lambda_*(E)=\lambda^*(E)
\]
を得る．
\end{proof}

\begin{lemma}[有限外測度での集合演算]\label{lem:finite-outer-measure-set-operations}
$E,F\subset\RR^d$ はLebesgue可測であり，
\[
\lambda^*(E)<\infty,\qquad \lambda^*(F)<\infty
\]
を満たすとする．
このとき
\[
E\cup F,\qquad E\setminus F,\qquad E\cap F
\]
も内外一致する．
さらに，$\{E_n\}_{n=1}^{\infty}$ がLebesgue可測集合列で，各 $n$ について
\[
\lambda^*(E_n)<\infty
\]
を満たし，
\[
B:=\bigcup_{n=1}^{\infty}E_n
\]
が $\lambda^*(B)<\infty$ を満たすならば，$B$ も内外一致する．
\end{lemma}
\begin{proof}
Lebesgue可測かつ外測度有限な集合は，
有限外測度集合の可測性（\cref{thm:finite-outer-measurable-iff-inner-outer}）より内外一致する．
したがって，$E,F$ および各 $E_n$ は内外一致する．
以下では\red{有限外測度集合の近似判定法（\cref{lem:finite-outer-approx-criterion}）の (3) }を用いる．

まず和集合を考える．
$\eps>0$ を任意に取る．
開集合 $G_E\supset E$ と $G_F\supset F$ を
\[
\lambda^*(G_E\setminus E)<\frac{\eps}{2},
\qquad
\lambda^*(G_F\setminus F)<\frac{\eps}{2}
\]
となるように取る．
すると $G:=G_E\cup G_F$ は $E\cup F$ を含む開集合であり，
\[
G\setminus(E\cup F)
\subset
(G_E\setminus E)\cup(G_F\setminus F)
\]
だから
\[
\lambda^*(G\setminus(E\cup F))<\eps.
\]
よって $E\cup F$ は内外一致する．

次に差集合を考える．
開集合 $G_E\supset E$ と，コンパクト集合 $K_F\subset F$ を
\[
\lambda^*(G_E\setminus E)<\frac{\eps}{2},
\qquad
\lambda^*(F\setminus K_F)<\frac{\eps}{2}
\]
となるように取る．
後者は\red{有限外測度集合の近似判定法（\cref{lem:finite-outer-approx-criterion}）の (2) }から従う．
このとき
\[
H:=G_E\setminus K_F
\]
は $E\setminus F$ を含む開集合である．
$x\in H\setminus(E\setminus F)$ ならば，$x\notin E$ または $x\in F$ である．
さらに $x\notin K_F$ だから
\[
H\setminus(E\setminus F)
\subset
(G_E\setminus E)\cup(F\setminus K_F).
\]
したがって
\[
\lambda^*(H\setminus(E\setminus F))<\eps.
\]
よって $E\setminus F$ も内外一致する．

共通部分は
\[
E\cap F=E\setminus(E\setminus F)
\]
と書けるので，すでに示した差集合の場合から従う．

最後に可算和を考える．
$\eps>0$ を任意に取る．
各 $n$ について開集合 $G_n\supset E_n$ を
\[
\lambda^*(G_n\setminus E_n)<\frac{\eps}{2^n}
\]
となるように取る．
このとき
\[
G:=\bigcup_{n=1}^{\infty}G_n
\]
は $B$ を含む開集合であり，
\[
G\setminus B
\subset
\bigcup_{n=1}^{\infty}(G_n\setminus E_n)
\]
だから，Lebesgue外測度の可算劣加法性（\cref{thm:lebesgue-outer-measure}）より
\[
\lambda^*(G\setminus B)
\le
\sum_{n=1}^{\infty}\lambda^*(G_n\setminus E_n)
<\eps.
\]
$\lambda^*(B)<\infty$ という仮定と\red{有限外測度集合の近似判定法
（\cref{lem:finite-outer-approx-criterion}）より}，$B$ は内外一致する．
\end{proof}

\begin{theorem}[Lebesgue可測集合の $\sigma$-加法族性]\label{thm:lebesgue-measurable-sigma-algebra}
Lebesgue可測集合全体 $\calL_d$ は $\RR^d$ 上の $\sigma$-加法族である．
すなわち，次が成り立つ．
\begin{enumerate}
    \item $\emptyset,\RR^d\in\calL_d$
    \item $A\in\calL_d$ ならば $\RR^d\setminus A\in\calL_d$
    \item $\{A_n\}_{n=1}^{\infty}\subset\calL_d$ ならば
    \[
    \bigcup_{n=1}^{\infty}A_n\in\calL_d
    \]
\end{enumerate}
\end{theorem}
\begin{proof}
$\emptyset\cap Q_R=\emptyset$ であり，また
$\RR^d\cap Q_R=Q_R$ はコンパクト集合である．
コンパクト集合では内外が一致すること（\cref{prop:compact-inner-outer-equal}）より，
\[
\emptyset,\RR^d\in\calL_d
\]
である．

$A\in\calL_d$ とする．
任意の $R>0$ に対して
\[
(\RR^d\setminus A)\cap Q_R
=
Q_R\setminus(A\cap Q_R)
\]
である．
$Q_R$ はコンパクトなのでLebesgue可測で外測度有限である．
また $A\cap Q_R$ は有界であり，定義より内外一致するので，
%\red{有界集合の可測性（\cref{lem:bounded-measurable-iff-inner-outer}）より}
Lebesgue可測である．
したがって有限外測度での集合演算（\cref{lem:finite-outer-measure-set-operations}）を適用でき，右辺は内外一致する．
したがって $\RR^d\setminus A\in\calL_d$ である．

次に $\{A_n\}_{n=1}^{\infty}\subset\calL_d$ とする．
任意の $R>0$ に対して
\[
\left(\bigcup_{n=1}^{\infty}A_n\right)\cap Q_R
=
\bigcup_{n=1}^{\infty}(A_n\cap Q_R)
\]
である．
各 $A_n\cap Q_R$ は有界であり，定義より内外一致するので，
%\red{有界集合の可測性（\cref{lem:bounded-measurable-iff-inner-outer}）より}
Lebesgue可測である．
また右辺は $Q_R$ に含まれるので有限外測度をもつ．
有限外測度での集合演算（\cref{lem:finite-outer-measure-set-operations}）の可算和の場合より，右辺は内外一致する．
ゆえに
\[
\bigcup_{n=1}^{\infty}A_n\in\calL_d
\]
である．
\end{proof}

\begin{corollary}[集合演算]\label{cor:lebesgue-measurable-set-operations}
$A,B\in\calL_d$ ならば
\[
A\cup B,\qquad A\cap B,\qquad A\setminus B
\]
はいずれもLebesgue可測である．
また $\{A_n\}_{n=1}^{\infty}\subset\calL_d$ ならば
\[
\bigcap_{n=1}^{\infty}A_n
\]
もLebesgue可測である．
\end{corollary}
\begin{proof}
Lebesgue可測集合の $\sigma$-加法族性（\cref{thm:lebesgue-measurable-sigma-algebra}）から従う．
例えば
\[
A\cap B=(A^c\cup B^c)^c,
\qquad
A\setminus B=A\cap B^c
\]
であり，可算共通部分は
\[
\bigcap_{n=1}^{\infty}A_n
=
\left(\bigcup_{n=1}^{\infty}A_n^c\right)^c
\]
と書ける．
\end{proof}

\begin{theorem}[開集合と閉集合]\label{thm:open-closed-lebesgue-measurable}
開集合と閉集合はLebesgue可測である．
\end{theorem}
\begin{proof}
まず開集合 $G\subset\RR^d$ を取る．
任意の $R>0$ に対して
\[
E:=G\cap Q_R
\]
とおく．
これは $Q_R$ に含まれるので有限外測度をもつ．
$\eps>0$ を任意に取る．
$\eta>0$ を十分小さく取って
\[
\lambda^*\bigl((-R-\eta,R+\eta)^d\setminus Q_R\bigr)<\eps
\]
となるようにする．
実際，左辺の集合はJordan可測であり，そのJordan測度は $\eta\downarrow0$ で $0$ に近づく．
このとき
\[
H:=G\cap(-R-\eta,R+\eta)^d
\]
は $E$ を含む開集合であり，
\[
H\setminus E
\subset
(-R-\eta,R+\eta)^d\setminus Q_R
\]
だから
\[
\lambda^*(H\setminus E)<\eps.
\]
\red{有限外測度集合の近似判定法（\cref{lem:finite-outer-approx-criterion}）より，}
$E=G\cap Q_R$ は内外一致する．
$R>0$ は任意なので $G$ はLebesgue可測である．

閉集合は開集合の補集合であり，Lebesgue可測集合の $\sigma$-加法族性
（\cref{thm:lebesgue-measurable-sigma-algebra}）より $\calL_d$ は補集合で閉じている．
したがって閉集合もLebesgue可測である．
\end{proof}

\section{Lebesgue測度の性質}

\begin{lemma}[有限外測度での有限加法性]\label{lem:finite-outer-measure-finite-additivity}
$E,F\subset\RR^d$ が有限外測度をもち内外一致しているとする．
さらに $E\cap F=\emptyset$ ならば
\[
\lambda^*(E\cup F)=\lambda^*(E)+\lambda^*(F)
\]
である．
\end{lemma}
\begin{proof}
外測度の劣加法性より
\[
\lambda^*(E\cup F)\le \lambda^*(E)+\lambda^*(F)
\]
である．
逆向きの不等式を示す．
$\eps>0$ を任意に取る．
内測度の定義より，コンパクト集合 $K_E\subset E$ と $K_F\subset F$ を
\[
\lambda^*(K_E)>\lambda^*(E)-\eps,
\qquad
\lambda^*(K_F)>\lambda^*(F)-\eps
\]
となるように取る．
$K_E$ と $K_F$ は互いに素なコンパクト集合だから，空でない場合は正の距離だけ離れている．
正の距離だけ離れた集合の加法性（\cref{lem:positive-distance-additive}）より
\[
\lambda^*(K_E\cup K_F)=\lambda^*(K_E)+\lambda^*(K_F)
\]
である．
したがってLebesgue外測度の単調性
（\cref{thm:lebesgue-outer-measure}）より
\[
\lambda^*(E\cup F)
\ge
\lambda^*(K_E\cup K_F)
>
\lambda^*(E)+\lambda^*(F)-2\eps.
\]
$\eps>0$ は任意なので
\[
\lambda^*(E\cup F)\ge \lambda^*(E)+\lambda^*(F)
\]
を得る．
\end{proof}

\begin{lemma}[有界窓による近似]\label{lem:bounded-window-approximation}
$A\in\calL_d$ に対して
\[
\lambda(A)=\lim_{R\to\infty}\lambda(A\cap Q_R)
=\sup_{R>0}\lambda(A\cap Q_R)
\]
が成り立つ．
\end{lemma}
\begin{proof}
$R<S$ なら $A\cap Q_R\subset A\cap Q_S$ であるから，
$\lambda(A\cap Q_R)$ は $R$ について単調増加である．
またLebesgue外測度の単調性（\cref{thm:lebesgue-outer-measure}）より
\[
\sup_{R>0}\lambda(A\cap Q_R)\le \lambda(A)
\]
である．

まず $\lambda(A)<\infty$ の場合を考える．
$\eps>0$ を任意に取る．
Lebesgue外測度の定義より，$A$ を覆う区間列 $\{I_n\}_{n=1}^{\infty}$ で
\[
\sum_{n=1}^{\infty}|I_n|<\lambda(A)+\eps
\]
となるものが取れる．
有限個の項を十分多く取れば，ある $N$ について
\[
\sum_{n>N}|I_n|<\eps
\]
である．
このとき $I_1,\dots,I_N$ は，空集合を除けばすべて有界である．
$I_1,\dots,I_N$ を含むように $R$ を十分大きく取ると，
\[
A\setminus Q_R\subset\bigcup_{n>N}I_n
\]
となるので
\[
\lambda^*(A\setminus Q_R)<\eps.
\]
外測度の劣加法性より
\[
\lambda(A)
\le
\lambda(A\cap Q_R)+\lambda^*(A\setminus Q_R)
<
\lambda(A\cap Q_R)+\eps.
\]
したがって
\[
\lambda(A)\le \sup_{R>0}\lambda(A\cap Q_R)
\]
である．

次に $\lambda(A)=\infty$ の場合を考える．
もし
\[
\sup_{R>0}\lambda(A\cap Q_R)<\infty
\]
なら，整数 $m\ge1$ に対して
\[
B_1:=A\cap Q_1,
\qquad
B_m:=A\cap(Q_m\setminus Q_{m-1})\quad(m\ge2)
\]
とおくと，$B_m$ たちは互いに素なLebesgue可測集合であり，しかも有界である．
有限外測度での有限加法性（\cref{lem:finite-outer-measure-finite-additivity}）を繰り返すと
\[
\sum_{m=1}^{N}\lambda(B_m)
=
\lambda(A\cap Q_N)
\le
\sup_{R>0}\lambda(A\cap Q_R)
\]
である．
よって
\[
\sum_{m=1}^{\infty}\lambda(B_m)<\infty.
\]
一方
\[
A=\bigcup_{m=1}^{\infty}B_m
\]
だから，Lebesgue外測度の可算劣加法性（\cref{thm:lebesgue-outer-measure}）より
\[
\lambda(A)
\le
\sum_{m=1}^{\infty}\lambda(B_m)
<\infty
\]
となり，$\lambda(A)=\infty$ に反する．
したがって
\[
\sup_{R>0}\lambda(A\cap Q_R)=\infty=\lambda(A)
\]
である．
\end{proof}

\begin{theorem}[完全加法性]\label{thm:lebesgue-measure-countable-additivity}
互いに素なLebesgue可測集合列 $\{A_n\}_{n=1}^{\infty}$ に対して
\[
\lambda\left(\bigsqcup_{n=1}^{\infty}A_n\right)
=
\sum_{n=1}^{\infty}\lambda(A_n)
\]
が成り立つ．
\end{theorem}
\begin{proof}
まず有限加法性を示す．
$A,B\in\calL_d$ が互いに素であるとする．
任意の $R>0$ に対して
\[
(A\cup B)\cap Q_R
=
(A\cap Q_R)\sqcup(B\cap Q_R)
\]
である．
右辺の2つの集合は有限外測度をもち内外一致しているので，
有限外測度での有限加法性（\cref{lem:finite-outer-measure-finite-additivity}）より
\[
\lambda((A\cup B)\cap Q_R)
=
\lambda(A\cap Q_R)+\lambda(B\cap Q_R)
\]
である．
$R\to\infty$ とし，有界窓による近似（\cref{lem:bounded-window-approximation}）を用いると
\[
\lambda(A\cup B)=\lambda(A)+\lambda(B)
\]
を得る．
帰納法により，互いに素な有限個のLebesgue可測集合についても有限加法性が成り立つ．

可算個の場合を示す．
Lebesgue外測度の可算劣加法性（\cref{thm:lebesgue-outer-measure}）より
\[
\lambda\left(\bigsqcup_{n=1}^{\infty}A_n\right)
\le
\sum_{n=1}^{\infty}\lambda(A_n)
\]
である．
逆向きについては，任意の $N$ に対して
Lebesgue外測度の単調性（\cref{thm:lebesgue-outer-measure}）と，
いま示した有限加法性より
\[
\lambda\left(\bigsqcup_{n=1}^{\infty}A_n\right)
\ge
\lambda\left(\bigsqcup_{n=1}^{N}A_n\right)
=
\sum_{n=1}^{N}\lambda(A_n).
\]
$N\to\infty$ とすれば
\[
\lambda\left(\bigsqcup_{n=1}^{\infty}A_n\right)
\ge
\sum_{n=1}^{\infty}\lambda(A_n)
\]
を得る．
\end{proof}

\begin{theorem}[基本性質]\label{thm:lebesgue-measure-basic}
Lebesgue測度について次が成り立つ．
\begin{enumerate}
    \item 非負性：
    \[
    0\le \lambda(A)\le\infty
    \qquad(A\in\calL_d)
    \]
    かつ $\lambda(\emptyset)=0$ である．
    \item 単調性：$A,B\in\calL_d$ かつ $A\subset B$ ならば
    \[
    \lambda(A)\le \lambda(B)
    \]
    \item 劣加法性：$\{A_n\}_{n=1}^{\infty}\subset\calL_d$ ならば
    \[
    \lambda\left(\bigcup_{n=1}^{\infty}A_n\right)
    \le
    \sum_{n=1}^{\infty}\lambda(A_n)
    \]
\end{enumerate}
\end{theorem}
\begin{proof}
\begin{enumerate}
    \item Lebesgue測度はLebesgue外測度の制限なので非負である．
    また $\lambda(\emptyset)=\lambda^*(\emptyset)=0$ である．

    \item Lebesgue外測度の単調性（\cref{thm:lebesgue-outer-measure}）から直ちに従う．

    \item
    \[
    B_1:=A_1,\qquad
    B_n:=A_n\setminus\bigcup_{k=1}^{n-1}A_k\quad(n\ge2)
    \]
    とおく．
    Lebesgue可測集合の集合演算（\cref{cor:lebesgue-measurable-set-operations}）より
    各 $B_n$ はLebesgue可測で，互いに素であり，
    \[
    \bigcup_{n=1}^{\infty}A_n=\bigsqcup_{n=1}^{\infty}B_n,
    \qquad
    B_n\subset A_n
    \]
    である．
    完全加法性（\cref{thm:lebesgue-measure-countable-additivity}）と
    上の単調性より
    \[
    \lambda\left(\bigcup_{n=1}^{\infty}A_n\right)
    =
    \sum_{n=1}^{\infty}\lambda(B_n)
    \le
    \sum_{n=1}^{\infty}\lambda(A_n)
    \]
    を得る．
\end{enumerate}
\end{proof}

\begin{theorem}[下からの連続性]\label{thm:lebesgue-measure-continuity-from-below}
Lebesgue可測集合の単調増大列
\[
A_1\subset A_2\subset\cdots
\]
に対して
\[
\lambda\left(\bigcup_{n=1}^{\infty}A_n\right)
=
\lim_{n\to\infty}\lambda(A_n)
\]
が成り立つ．
\end{theorem}
\begin{proof}
\[
B_1:=A_1,\qquad
B_n:=A_n\setminus A_{n-1}\quad(n\ge2)
\]
とおく．
すると $B_n$ たちは互いに素であり，
\[
A_N=\bigsqcup_{n=1}^{N}B_n,
\qquad
\bigcup_{n=1}^{\infty}A_n=\bigsqcup_{n=1}^{\infty}B_n
\]
である．
完全加法性（\cref{thm:lebesgue-measure-countable-additivity}）より
\[
\lambda(A_N)=\sum_{n=1}^{N}\lambda(B_n),
\qquad
\lambda\left(\bigcup_{n=1}^{\infty}A_n\right)
=
\sum_{n=1}^{\infty}\lambda(B_n).
\]
$N\to\infty$ とすれば主張を得る．
\end{proof}

\begin{theorem}[上からの連続性]\label{thm:lebesgue-measure-continuity-from-above}
Lebesgue可測集合の単調減少列
\[
A_1\supset A_2\supset\cdots
\]
に対して，$\lambda(A_1)<\infty$ ならば
\[
\lambda\left(\bigcap_{n=1}^{\infty}A_n\right)
=
\lim_{n\to\infty}\lambda(A_n)
\]
が成り立つ．
\end{theorem}
\begin{proof}
\[
B_n:=A_1\setminus A_n
\]
とおく．
すると
\[
B_1\subset B_2\subset\cdots,
\qquad
\bigcup_{n=1}^{\infty}B_n
=
A_1\setminus\bigcap_{n=1}^{\infty}A_n
\]
である．
下からの連続性（\cref{thm:lebesgue-measure-continuity-from-below}）より
\[
\lambda\left(A_1\setminus\bigcap_{n=1}^{\infty}A_n\right)
=
\lim_{n\to\infty}\lambda(A_1\setminus A_n).
\]
$A_n\subset A_1$ かつ $\lambda(A_1)<\infty$ なので，
完全加法性（\cref{thm:lebesgue-measure-countable-additivity}）から従う有限加法性より
\[
\lambda(A_1\setminus A_n)=\lambda(A_1)-\lambda(A_n)
\]
である．
同様に
\[
\lambda\left(A_1\setminus\bigcap_{n=1}^{\infty}A_n\right)
=
\lambda(A_1)-\lambda\left(\bigcap_{n=1}^{\infty}A_n\right).
\]
これらを合わせれば主張が従う．
\end{proof}

\begin{proposition}[平行移動不変性]\label{prop:lebesgue-measure-translation-invariant}
$A\in\calL_d$ と $c\in\RR^d$ に対して
\[
A+c:=\{x+c\mid x\in A\}
\]
もLebesgue可測であり，
\[
\lambda(A+c)=\lambda(A)
\]
が成り立つ．
\end{proposition}
\begin{proof}
まず可測性を示す．
任意の $R>0$ に対して
\[
E:=(A+c)\cap Q_R
\]
とおく．
平行移動して戻すと
\[
E-c=A\cap(Q_R-c)
\]
である．
$Q_R-c$ はコンパクト集合なのでLebesgue可測であり，
Lebesgue可測集合の集合演算（\cref{cor:lebesgue-measurable-set-operations}）より
$A\cap(Q_R-c)$ もLebesgue可測である．
これは有界集合だから，内外一致する．

Lebesgue外測度の平行移動不変性（\cref{prop:lebesgue-outer-translation-invariant}）を用いる．
またコンパクト集合の平行移動はコンパクト集合なので，内測度も平行移動で不変である．
したがって
\[
\lambda^*(E)=\lambda^*(E-c),
\qquad
\lambda_*(E)=\lambda_*(E-c)
\]
である．
$E-c$ は内外一致するので $E$ も内外一致する．
$R>0$ は任意だから $A+c$ はLebesgue可測である．

最後に，Lebesgue外測度の平行移動不変性（\cref{prop:lebesgue-outer-translation-invariant}）より
\[
\lambda(A+c)=\lambda^*(A+c)=\lambda^*(A)=\lambda(A)
\]
である．
\end{proof}
